%% LyX 2.4.4 created this file.  For more info, see https://www.lyx.org/.
%% Do not edit unless you really know what you are doing.
\documentclass[english]{article}
\usepackage[T1]{fontenc}
\usepackage[latin9]{inputenc}
\usepackage{units}
\usepackage{enumitem}
\usepackage{amsmath}
\usepackage{amsthm}
\usepackage{geometry}
\geometry{verbose,tmargin=1in,bmargin=1in,lmargin=1in,rmargin=1in}

\makeatletter
%%%%%%%%%%%%%%%%%%%%%%%%%%%%%% Textclass specific LaTeX commands.
\numberwithin{equation}{section}
\numberwithin{figure}{section}
\newlength{\lyxlabelwidth}      % auxiliary length 

%%%%%%%%%%%%%%%%%%%%%%%%%%%%%% User specified LaTeX commands.
\usepackage{fancyhdr}
\usepackage{lastpage}
\pagestyle{fancy}
\usepackage{enumitem}
\usepackage{ifthen}
%sets math fonts to be more readable
\everymath=\expandafter{\the\everymath\displaystyle}
%\DeclareMathSizes{10}{10}{10}{10}
\usepackage{multicol}
% Added by lyx2lyx
\setlength{\parskip}{\smallskipamount}
\setlength{\parindent}{0pt}

\makeatother

\usepackage{babel}
\begin{document}
\renewcommand{\author}{Dr.\,James Ashe}

\newcommand{\classNum}{231}

\newcommand{\classSec}{C}

\newcommand{\nameType}{Name:}

\newcommand{\examType}{Exam 2 Review/Quiz}

\renewcommand{\date}{Due November 11, 2012}

%%First page's header%%%%%%%%%%%%%%%%%%%%%%
\fancypagestyle{firststyle} %%header settings for first pagr
{
	\fancyhead[L]{MTH \classNum{}(\classSec{}) \author{}\\
	\textbf{\examType{}} \date{}}
	\fancyhead[C]{\textbf{\nameType}}
	\setlength{\headsep}{14pt}
}
%%%%%%%%%%%%%%%%%%%%%%%%%%%%%%%%%%%%%%%%%%

%%Next page's header%%%%%%%%%%%%%%%%%%%%%%%
\setlist{nolistsep}
\lhead{\classNum{}(\classSec{})} 
\chead{\textbf{\nameType}} 
\cfoot{\thepage{}~of~\pageref{LastPage}} 
\setlength{\headsep}{5pt} 
%%%%%%%%%%%%%%%%%%%%%%%%%%%%%%%%%%%%%%%%%%%

%%Various settings; see the preamble for other settings%%%%%
%%\setenumerate[0]{leftmargin=12pt,itemindent=0pt,labelwidth=0pt}
\renewcommand{\labelenumi}{\textbf{\arabic{enumi}.}}
\renewcommand{\labelenumii}{\textbf{\alph{enumii}.}}
\thispagestyle{firststyle}
%%%%%%%%%%%%%%%%%%%%%%%%%%%%%%%%%%%%%%%%%%

\section*{6.1: Velocity and Net Change}

\textbf{HW problems 7-11 (odd). }\emph{Find the displacement, and
find the distance traveled over the given interval.}
\begin{enumerate}
\item $v(t)=10\sin2t$; $0\leq t\leq2\pi$\vfill
\end{enumerate}
\textbf{HW problems 12 and 15. }\emph{Determine the position function
(you may use EITHER the FTC or anti-derivative method), and determine
when the object is moving in the positive and negative directions.}
\begin{enumerate}[resume]
\item $v(t)=\frac{20}{\sqrt{t+1}}$ on $[0,8]$; $s(0)=-4$ (use substitution)\vfill
\end{enumerate}
\textbf{HW problems 18, 25, 26, 31, and 32. }\emph{Word problems applying
the same concepts as above.}
\begin{enumerate}[resume]
\item A culture of bacteria in a petri dish has an initial population of
1500 cells andgrows at a rate (in cells/day) of $N'(t)=200(t+2)^{\nicefrac{-1}{2}}$.

\begin{enumerate}[resume]
\item What is the population after 14 days? 
\item Find the population $N(t)$ at any time $t\geq0$.\vfill
\end{enumerate}
\end{enumerate}

\section*{6.2 Region Between Curves}

\textbf{HW problems 5-6, 10-11, 15-18, 23-24 }\emph{Find the area
of the region.}
\begin{enumerate}[resume]
\item The region bounded by $y=x$ and $y=x^{2}-2$.\vfill
\item The region in the first quadrant bounded by $y=(x-1)^{3}$ and $y=x-1$.\vfill
\item The region bounded by $y=x$ and $x=(y-2)^{2}$.\vfill
\end{enumerate}

\section*{6.3 Volume by Slicing}

\textbf{HW problems 7 and 9. }\emph{Use the general slicing method
to find the volume of the following solids.}
\begin{enumerate}[resume]
\item A triangular wedge whose perpendicular sides have lengths 3, 4, and
5 (use calculus; see the figure on page 343).\vfill
\item The solid with a semicircular base of radius 5 whose cross sections
are perpendicular to the base and parallel to the diameter. \vfill
\end{enumerate}
\textbf{HW problems 15-18. }\emph{Let R be the region bounded by the
following curves. Use the disk method to find the volume of the solid
generated when $R$ is revolved around the $x-$axis.}
\begin{enumerate}[resume]
\item $y=2x$, $y=0$, $x=3$ (see the figure on page 344).\vfill
\item $y=4-x^{2}$, $y=0$, $x=0$ (see the figure on page 344). \vfill
\item $y=\cos x$, $y=0$, $x=0$ (Recall that $\cos^{2}x=\frac{1}{2}\left(1+\cos2x\right)$.)\vfill
\end{enumerate}

\section*{6.5 Length of Curves}

\textbf{HW problems 3-9 odd. }\emph{Find the arc length of the following
curves on the given interval by integrating with respect to $x$.}
\begin{enumerate}[resume]
\item $y=\frac{1}{3}x^{\nicefrac{3}{2}}$; $[0,60]$\vfill
\item $y=\frac{(x^{2}+2)^{\nicefrac{3}{2}}}{3}$; $[0,1]$\vfill
\item $y=\frac{x^{4}}{4}+\frac{1}{8x^{2}}$; $[1,2]$\vfill
\end{enumerate}

\end{document}
